\section{Introduction}

Dynamic systems with springs and dampeners are used widely in engineering applications. Spring, dampener systems are mainly used to react to and suppress external forces acting on a system. An example of this is a car's suspension system, which is used to minimise the vibrations of the chassis when on bumpy or rough surfaces. This report is an investigation into car suspension systems. Modelling these dynamic systems is crucial to understand their behaviour under different conditions, which can then be used to optimise performance, stability and reliability. We aim to determine the effectiveness of different suspension systems by finding a meaningful correlation between dampening and spring coefficients, and the safety of the passengers.

To do this we have made a computational model of our system. Our model represents the main body of the car as a beam suspended by two different spring and dampener systems on either side. This system is mechanically interesting as it extends the dampened spring and pendulum systems seen throughout the semester with further ridged body dynamics.

Our report contains a detailed methodology that describes in detail the mathematical derivation of our equations of motion from the base system. These equations of motion are the basis that the computational method will be founded on. Our code is broken up into three major sections; definition of variables and equations, solving of equations, and the extraction of useful data, later shown through diagrams. Finally, our findings are discussed at length in the results section.

Throughout this investigation in ridged body dynamics we hope to better understand suspension, the vital safety mechanisms within all our cars.
